% !TeX program = xelatex

% Author: Christopher Métrailler
% Date: June 26, 2016
% Version: v0.1
% 
% XeLaTeX is required to build this template file.

\documentclass{article}
\usepackage{fontspec}

\newlength{\leftcolwidth}
\newlength{\rightcolwidth}

\usepackage[hidelinks]{hyperref}

\usepackage{geometry}
\geometry{a4paper, landscape} % 297mm x 210mm (landscape)

% \usepackage{xcolor}
\usepackage[table]{xcolor}

\usepackage{graphicx}
\graphicspath{ {images/} }

\usepackage{tikz}
\usetikzlibrary{calc}

\parindent=0mm

% http://tex.stackexchange.com/questions/25249/how-do-i-use-a-particular-font-for-a-small-section-of-text-in-my-document
% http://tex.stackexchange.com/questions/62603/how-to-choose-a-specific-weight-from-a-font-family-using-fontspec-and-xelatex

% Custom Oswald Regular / Light fonts
\newfontfamily{\oswald}[
  Path        = fonts/,
  Ligatures   = TeX,
  Extension   = .ttf,  
  UprightFont = *-Regular,
  ItalicFont  = *-Light
]{Oswald}

% Oswald Regular font
\newcommand{\orfont}[2]{\fontsize{#1}{#1}\oswald{\color{white}{#2}}}

% Oswald Light font
\newcommand{\olfont}[2]{\fontsize{#1}{#1}\oswald{\color{white}\textit{#2}}}


% Custom Muli Light font
\newfontfamily{\muli}[
Path        = fonts/,
Ligatures   = TeX,
Extension   = .ttf,  
ItalicFont  = *-Light
]{Muli}

% Muli Light font
\newcommand{\mlfont}[2]{\fontsize{#1}{#1}\muli{\color{white}\textit{#2}}}



\newenvironment{nscenter}
{\parskip=0pt\par\nopagebreak\centering}
{\par\noindent\ignorespacesafterend}

\begin{document}
\pagestyle{empty}

\setlength{\leftcolwidth}{106mm}           % Should corresponds to the background image width
\setlength{\rightcolwidth}{\paperwidth}
\addtolength\rightcolwidth{-\leftcolwidth} % Fit the width of the paper

% \fontsize{size}{skip}

% Custom column style
\tikzstyle{col} = [anchor=south west, rectangle, inner sep=0pt, inner ysep=0pt, minimum width=#1, minimum height=\paperheight]

% From http://tex.stackexchange.com/questions/260775/how-to-place-text-over-an-image
% and http://tex.stackexchange.com/questions/296659/remove-tikzpicture-border

% From http://tex.stackexchange.com/questions/199968/how-do-i-set-only-one-side-of-a-tikz-nodes-border-color

% Left column
% -----------

\begin{tikzpicture}[remember picture, overlay]

% Background image of the column
\node[col=\leftcolwidth] (X) at (current page.south west) {
  \includegraphics[height=\paperheight]{col_left}
};

\begin{scope}[x={(X.south west)}, y={(X.north east)}]%

  % Column background color
  \fill[fill=blue, fill opacity=0.7] ($(X.south west)$) rectangle ($(X.north east)$);


  \node[fill=blue, text width=0.8\leftcolwidth, anchor=north, yshift=-10mm] (A) at (X.north) {
  	\centering \olfont{14}{Open-Source} \\
  };

  \node[fill=red, text width=0.8\leftcolwidth, anchor=center, yshift=-5mm] (B) at (A.south) {
    \centering \orfont{16}{BLUETOOTH 4.2} \olfont{16}{SENSOR BEACON} \\
  };

  \node[fill=red, text width=0.8\leftcolwidth, anchor=center, yshift=-20mm] (C) at (B.south) {
    \centering \orfont{32}{RuuviTag} \\
  };

  \node[fill=red, anchor=center, yshift=-40mm] (D) at (C.south) {
    \includegraphics[height=30mm]{col_left}
  };

  % Build from the bottom

  \node[anchor=center, yshift=15mm] (G) at (X.south) {
	\centering \mlfont{11}{\href{http://ruuvitag.com}{ruuvitag.com}}
  };

  \node[fill=red, anchor=center, yshift=15mm] (F) at (G.north) {
	\includegraphics[height=15mm]{col_left}
  };

  \node[, anchor=center, yshift=30mm] (E) at (F.south) {
	\centering \orfont{14}{KICK}
  };

  \draw[ultra thick, yellow] ($(E.south west)+(4pt,0)$) |- ($(E.south east)-(4pt,0)$);

\end{scope}

\end{tikzpicture}

%\pagestyle{empty}%
%    \noindent
%    \begin{tabular}{@{}c@{}}
%    \cellcolor{blue!50}%
%    \begin{minipage}[t][\paperheight][t]{106mm}%
%        Left
%        \includegraphics[width=\linewidth]{col_left} 
%        sdsdf
%    \end{minipage}%
%    \end{tabular}%
%    \begin{tabular}{@{}c@{}}
%    %\cellcolor{green!50}%
%    \begin{minipage}[t][\paperheight][t]{191mm}%
%        Right
%    \end{minipage}%
%    \end{tabular}%


\end{document}